{\rtf1\ansi\ansicpg1252\cocoartf1504\cocoasubrtf830
{\fonttbl\f0\fswiss\fcharset0 Helvetica;}
{\colortbl;\red255\green255\blue255;}
{\*\expandedcolortbl;;}
\paperw11900\paperh16840\margl1440\margr1440\vieww10800\viewh8400\viewkind0
\pard\tx566\tx1133\tx1700\tx2267\tx2834\tx3401\tx3968\tx4535\tx5102\tx5669\tx6236\tx6803\pardirnatural\partightenfactor0

\f0\fs24 \cf0 \\documentclass[9pt,french]\{book\}\
\\usepackage\{teach\}\
\\usepackage\{verbatim\}\
\\usepackage\{amsmath,amsfonts,amssymb\}\
\\usepackage\{pgfplots\}\
\\usepackage\{eurosym\}\
\\usepackage\{geometry\}\
\\geometry\{top=1cm, bottom=1cm, left=2cm, right=2cm\}\
\\usepackage[utf8]\{inputenc\}\
\\usepackage[french]\{babel\}\
\\redefinecolor\{thm\}\{title\}\{bgcolor\}\{alizarin\}\
\\redefinecolor\{prop\}\{title\}\{bgcolor\}\{meatbrown\}\
\\redefinecolor\{defn\}\{title\}\{bgcolor\}\{olivine\}\
\\definecolor\{alizarin\}\{rgb\}\{0.82, 0.1, 0.26\}\
\\definecolor\{meatbrown\}\{rgb\}\{0.9, 0.72, 0.23\}\
\\definecolor\{olivine\}\{rgb\}\{0.6, 0.73, 0.45\}\
\\usepackage\{pgf,tikz,tkz-tab\}\
\
\\usepackage\{mathrsfs\}\
\\usetikzlibrary\{arrows\}\
\
\\begin\{document\}\
\\twocolumn\
\\underline\{Devoir \'e0 la maison\}\
\
\\underline\{Exercice 1\}\
\\begin\{itemize\}\
\
\\item[1)] Dans une classe de 30 \'e9l\'e8ves, il y \'e0 40\\% de sportifs r\'e9guliers. Quel est le nombre de sportifs r\'e9guliers dans la classe\'a0?\
\
\\item[2)] 20\\% des 1000 \'e9l\'e8ves d\'92un lyc\'e9e sont externes. \
Combien y-a-t\'92il d\'92\'e9l\'e8ves externes au lyc\'e9e\'a0?\
\
\\end\{itemize\}\
\
\\underline\{Exercice 2\}\
\
Dire si les propositions sont vraies ou fausses et justifier la r\'e9ponse.\
\
\\begin\{itemize\}\
\
\\item[1)]  Augmenter une quantit\'e9 de 100\\% revient \'e0 la multiplier par 2.\
\\item[2)] Pour baisser une quantit\'e9 de 15\\%, on la multiplie par 0,15.\
\\item[3)]  Pour augmenter une quantit\'e9 de 7\\%, on la multiplie par 1,07.\
\
\\end\{itemize\}\
\
\\underline\{Exercice 3\}\
\
Un magasin de jeans a achet\'e9 des pantalons au prix de gros de 25\\euro\{\} le pantalon. Il majore habituellement le prix de 40\\% pour les vendre.\
\
\\begin\{itemize\}\
    \\item[1)] Calculer la diff\'e9rence entre le prix de vente et le prix de gros. Il s\'92agit de la marge que fait le magasin.\
    \\item[2)] Le magasin d\'e9cide de solder les derniers jeans qui lui restent. Il applique une r\'e9duction de 25\\% au moment des f\'eates de No\'ebl. A quel prix va-t-il vendre ses jeans\'a0? \\\\\
    \
\\end\{itemize\}\
\
\
\\underline\{Devoir \'e0 la maison\}\
\
\\underline\{Exercice 1\}\
\\begin\{itemize\}\
\
\\item[1)] Dans une classe de 30 \'e9l\'e8ves, il y \'e0 40\\% de sportifs r\'e9guliers. Quel est le nombre de sportifs r\'e9guliers dans la classe\'a0?\
\
\\item[2)] 20\\% des 1000 \'e9l\'e8ves d\'92un lyc\'e9e sont externes. \
Combien y-a-t\'92il d\'92\'e9l\'e8ves externes au lyc\'e9e\'a0?\
\
\\end\{itemize\}\
\
\\underline\{Exercice 2\}\
\
Dire si les propositions sont vraies ou fausses et justifier la r\'e9ponse.\
\
\\begin\{itemize\}\
\
\\item[1)]  Augmenter une quantit\'e9 de 100\\% revient \'e0 la multiplier par 2.\
\\item[2)] Pour baisser une quantit\'e9 de 15\\%, on la multiplie par 0,15.\
\\item[3)]  Pour augmenter une quantit\'e9 de 7\\%, on la multiplie par 1,07.\
\
\\end\{itemize\}\
\
\\underline\{Exercice 3\}\
\
Un magasin de jeans a achet\'e9 des pantalons au prix de gros de 25\\euro\{\} le pantalon. Il majore habituellement le prix de 40\\% pour les vendre.\
\
\\begin\{itemize\}\
    \\item[1)] Calculer la diff\'e9rence entre le prix de vente et le prix de gros. Il s\'92agit de la marge que fait le magasin.\
    \\item[2)] Le magasin d\'e9cide de solder les derniers jeans qui lui restent. Il applique une r\'e9duction de 25\\% au moment des f\'eates de No\'ebl. A quel prix va-t-il vendre ses jeans\'a0?\
\\end\{itemize\}\
\
\\underline\{Devoir \'e0 la maison\}\
\
\\underline\{Exercice 1\}\
\\begin\{itemize\}\
\
\\item[1)] Dans une classe de 30 \'e9l\'e8ves, il y \'e0 40\\% de sportifs r\'e9guliers. Quel est le nombre de sportifs r\'e9guliers dans la classe\'a0?\
\
\\item[2)] 20\\% des 1000 \'e9l\'e8ves d\'92un lyc\'e9e sont externes. \
Combien y-a-t\'92il d\'92\'e9l\'e8ves externes au lyc\'e9e\'a0?\
\
\\end\{itemize\}\
\
\\underline\{Exercice 2\}\
\
Dire si les propositions sont vraies ou fausses et justifier la r\'e9ponse.\
\
\\begin\{itemize\}\
\
\\item[1)]  Augmenter une quantit\'e9 de 100\\% revient \'e0 la multiplier par 2.\
\\item[2)] Pour baisser une quantit\'e9 de 15\\%, on la multiplie par 0,15.\
\\item[3)]  Pour augmenter une quantit\'e9 de 7\\%, on la multiplie par 1,07.\
\
\\end\{itemize\}\
\
\\underline\{Exercice 3\}\
\
Un magasin de jeans a achet\'e9 des pantalons au prix de gros de 25\\euro\{\} le pantalon. Il majore habituellement le prix de 40\\% pour les vendre.\
\
\\begin\{itemize\}\
    \\item[1)] Calculer la diff\'e9rence entre le prix de vente et le prix de gros. Il s\'92agit de la marge que fait le magasin.\
    \\item[2)] Le magasin d\'e9cide de solder les derniers jeans qui lui restent. Il applique une r\'e9duction de 25\\% au moment des f\'eates de No\'ebl. A quel prix va-t-il vendre ses jeans\'a0?\\\\\
    \
\\end\{itemize\}\
\
\
\
\
\\underline\{Devoir \'e0 la maison\}\
\
\\underline\{Exercice 1\}\
\\begin\{itemize\}\
\
\\item[1)] Dans une classe de 30 \'e9l\'e8ves, il y \'e0 40\\% de sportifs r\'e9guliers. Quel est le nombre de sportifs r\'e9guliers dans la classe\'a0?\
\
\\item[2)] 20\\% des 1000 \'e9l\'e8ves d\'92un lyc\'e9e sont externes. \
Combien y-a-t\'92il d\'92\'e9l\'e8ves externes au lyc\'e9e\'a0?\
\
\\end\{itemize\}\
\
\\underline\{Exercice 2\}\
\
Dire si les propositions sont vraies ou fausses et justifier la r\'e9ponse.\
\
\\begin\{itemize\}\
\
\\item[1)]  Augmenter une quantit\'e9 de 100\\% revient \'e0 la multiplier par 2.\
\\item[2)] Pour baisser une quantit\'e9 de 15\\%, on la multiplie par 0,15.\
\\item[3)]  Pour augmenter une quantit\'e9 de 7\\%, on la multiplie par 1,07.\
\
\\end\{itemize\}\
\
\\underline\{Exercice 3\}\
\
Un magasin de jeans a achet\'e9 des pantalons au prix de gros de 25\\euro\{\} le pantalon. Il majore habituellement le prix de 40\\% pour les vendre.\
\
\\begin\{itemize\}\
    \\item[1)] Calculer la diff\'e9rence entre le prix de vente et le prix de gros. Il s\'92agit de la marge que fait le magasin.\
    \\item[2)] Le magasin d\'e9cide de solder les derniers jeans qui lui restent. Il applique une r\'e9duction de 25\\% au moment des f\'eates de No\'ebl. A quel prix va-t-il vendre ses jeans\'a0?\
\\end\{itemize\}\
\
\
\
\
\
\
\
\
\
\\end\{document\}\
\
}